\chapter{Grupy Liego i ich algebry}
\label{ch:grupy-liego}
\index{grupa Liego}
\index{algebra Liego}

W Rozdziale~\ref{ch:rozmaitosci} grupa ortogonalna pojawiła się jako
rozmaitość, a w Rozdziale~\ref{ch:pola} poznaliśmy przepływy i nawias pól.
Teraz połączymy te konstrukcje. W grupie Liego każdy wektor styczny
przy elemencie neutralnym wyznacza pole na całej grupie, a przepływ
tego pola daje rodzinę symetrii zależną od jednego parametru.
Grupy i homomorfizmy wystąpiły już w
Definicji~\ref{def:grupa-hom-5}; tutaj wymagamy ponadto gładkości.

\section{Grupa z gładką strukturą}
\label{sec:grupy-liego-przyklady}

\begin{definicja}[Grupa Liego]
\label{def:grupa-liego}
\index{grupa Liego!definicja}
\emph{Grupą Liego} nazywamy skończenie wymiarową gładką rozmaitość
bez brzegu $G$, wyposażoną w strukturę grupy, dla której
\[
 m:G\times G\longrightarrow G,\quad m(g,h)=gh,
 \qquad \iota:G\longrightarrow G,\quad\iota(g)=g^{-1}
\]
są gładkie. Jej element neutralny oznaczamy $e$.
Intuicyjnie można w niej wykonywać działania grupowe i jednocześnie
różniczkować je względem parametrów. Homomorfizm grup Liego nazywamy
\emph{gładkim}, gdy jest gładki jako odwzorowanie rozmaitości.
\end{definicja}

Dla $g\in G$ oznaczamy przez $L_g(h)=gh$ i $R_g(h)=hg$
translacje. Obie są dyfeomorfizmami: są gładkie jako ograniczenia
$m$, a ich odwrotnościami są odpowiednio $L_{g^{-1}}$ i
$R_{g^{-1}}$. Pozwalają przenosić rachunek z punktu $e$ do
każdego punktu grupy.

\begin{przyklad}[Grupy macierzy odwracalnych]
\label{prz:gl-liego}
\index{grupa ogólna liniowa}
Grupa $\mathrm{GL}(n,\R)$ z
Dodatku~\ref{sec:app-grupa-gl} jest otwartym podzbiorem
$\R^{n^2}$: jest przeciwobrazem otwartego $\R_*$
przez ciągły wielomian $\det$.
Iloczyn macierzy jest wielomianowy w ich współczynnikach, a
\[
 A^{-1}=\frac{\operatorname{adj}A}{\det A}
\]
jest gładki na tym otwartym zbiorze. Stąd $\mathrm{GL}(n,\R)$
jest grupą Liego wymiaru $n^2$.
W użytym wzorze $\operatorname{adj}A$ jest transponowaną macierzą
dopełnień algebraicznych. Rozwinięcie Laplace'a daje
$A\,\operatorname{adj}A=(\det A)I$: na przekątnej jest to
rozwinięcie $\det A$, a poza nią wyznacznik macierzy
z dwoma jednakowymi wierszami, równy zeru.
Dzielenie przez niezerowy $\det A$ daje zatem odwrotność;
jej współczynniki są ilorazami wielomianów.
Ten sam argument, po potraktowaniu
części rzeczywistych i urojonych współczynników jako współrzędnych,
daje $\mathrm{GL}(n,\C)$ wymiaru rzeczywistego $2n^2$.
W obu przypadkach, dla $\mathbb K\in\{\R,\C\}$,
$T_I\mathrm{GL}(n,\mathbb K)=M_n(\mathbb K)$
jako przestrzeń rzeczywista: otwarty podzbiór przestrzeni macierzy
ma tę samą przestrzeń styczną co otoczenie.
\end{przyklad}

\begin{przyklad}[Grupy ortogonalne]
\label{prz:orth-liego}
\index{grupa ortogonalna}
W przykładzie „Grupa ortogonalna”, po
Twierdzeniu~\ref{tw:przeciwobraz}, wykazaliśmy, że
grupa $\mathrm O(n)$ z Dodatku~\ref{sec:app-grupa-o}
jest gładką rozmaitością i że
$T_I\mathrm O(n)=\{X:X^T=-X\}$. Mnożenie jest ograniczeniem
wielomianowego mnożenia macierzy, a odwrotność ma tutaj prostą
postać $A^{-1}=A^T$, więc oba działania są gładkie.
Ponieważ $\det A\in\{-1,1\}$ na $\mathrm O(n)$,
$\mathrm{SO}(n)$ jest
otwartą i domkniętą podrozmaitością $\mathrm O(n)$.
Istotnie, jest zarówno przeciwobrazem domkniętego $\{1\}$,
jak i otwartego $(0,\infty)$ przez ciągły wyznacznik.
Jest podgrupą Liego tego samego wymiaru $n(n-1)/2$.
Jej przestrzeń styczna przy $I$ również składa się z
macierzy skośnie symetrycznych.
\end{przyklad}

\begin{przyklad}[Torus jako grupa]
\label{prz:torus-liego}
\index{torus!jako grupa Liego}
Na torusie $T^k=\R^k/\Z^k$ z Przykładu~\ref{prz:torus}
kładziemy $[x]+[y]=[x+y]$ i $-[x]=[-x]$.
Zmiana reprezentantów na $x+m$, $y+n$, gdzie $m,n\in\Z^k$,
zmienia sumę o $m+n\in\Z^k$, a przeciwny wektor o $-m$.
Działania nie zależą więc od wyboru reprezentantów.
W lokalnych mapach
utworzonych z rzutowania $\R^k\to T^k$ działania mają postać
zwykłego dodawania i zmiany znaku, z ewentualną całkowitą
translacją współrzędnych. Są więc gładkie. Równoważny model
\[
 T^k\longrightarrow(S^1)^k,\qquad
 [x_1,\ldots,x_k]\longmapsto
 (e^{2\pi i x_1},\ldots,e^{2\pi i x_k})
\]
przekształca dodawanie w mnożenie zespolone na każdej współrzędnej;
odwrotność staje się sprzężeniem zespolonym.
Równość $e^{2\pi i x}=e^{2\pi i y}$ zachodzi dokładnie dla
$x-y\in\Z$, co daje iniektywność na klasach, a każdy punkt
okręgu ma postać $e^{2\pi i x}$, co daje surjektywność.
Na krótkim łuku wybieramy gładką gałąź kąta $\theta$;
lokalna odwrotność ma współrzędną $\theta/(2\pi)$.
Otrzymujemy zatem dyfeomorfizm iloczynów.
Dla $k=2$ wzór
$(u,v)\mapsto((R+r\cos v)\cos u,(R+r\cos v)\sin u,r\sin v)$,
gdzie $R>r>0$, wyznacza dyfeomorfizm
$(\R/2\pi\Z)^2$ na torus obrotowy z
Rozdziału~\ref{ch:rozmaitosci}. Przenosi on na ten torus
strukturę grupy. Jego mnożenie zależy od tej identyfikacji, a nie
od dodawania punktów w otaczającej go przestrzeni $\R^3$.
Lokalną odwrotność torusa obrotowego także można odczytać:
dla punktu $(x,y,z)$ kąt $u$ jest argumentem $x+iy$,
a $v$ argumentem $(\sqrt{x^2+y^2}-R)+iz$.
Obie liczby są niezerowe, gdyż ich moduły wynoszą odpowiednio
$R+r\cos v\geq R-r>0$ i $r>0$.
Lokalne gałęzie argumentu są gładkie; globalnie kąty
rozumiemy modulo $2\pi$.
\end{przyklad}

\begin{przyklad}[Macierze specjalnie unitarne]
\label{prz:su2-liego}
\index{grupa specjalna unitarna}
Niech $\mathrm{SU}(2)$ będzie grupą z
Dodatku~\ref{sec:app-grupa-u}.
Z warunków $U^{-1}=U^*$ i $\det U=1$ wynika, przez porównanie
wzorów na odwrotność, że każdy jej element ma jedyną postać
\begin{equation}\label{eq:su2-macierz}
 U(a,b)=\begin{pmatrix}a&b\\-\bar b&\bar a\end{pmatrix},
 \qquad |a|^2+|b|^2=1.
\end{equation}
Rozpiszmy porównanie. Dla $U=\begin{psmallmatrix}a&b\\c&d\end{psmallmatrix}$
o wyznaczniku $1$ wzór na odwrotność daje
$U^{-1}=\begin{psmallmatrix}d&-b\\-c&a\end{psmallmatrix}$,
a sprzężenie hermitowskie
$U^*=\begin{psmallmatrix}\bar a&\bar c\\\bar b&\bar d\end{psmallmatrix}$.
Porównanie współczynników daje $d=\bar a$, $c=-\bar b$.
Wyznacznik jest wtedy równy $|a|^2+|b|^2$, więc warunek
$\det U=1$ jest równaniem sfery.
Odwrotnie,
\[
 \begin{pmatrix}\bar a&-b\\\bar b&a\end{pmatrix}
 \begin{pmatrix}a&b\\-\bar b&\bar a\end{pmatrix}
 =
 \begin{pmatrix}|a|^2+|b|^2&\bar a b-b\bar a\\
                 \bar b a-a\bar b&|a|^2+|b|^2\end{pmatrix}=I.
\]
Współrzędne $(\Re a,\Im a,\Re b,\Im b)$ identyfikują zatem
$\mathrm{SU}(2)$ gładko ze sferą $S^3\subset\R^4$.
Mnożenie macierzy ogranicza się do gładkiego odwzorowania
$S^3\times S^3\to S^3$, a odwrotność $U^{-1}=U^*$ jest
ograniczeniem rzeczywiście liniowego odwzorowania na $M_2(\C)$.
Jest to więc zwarta, trójwymiarowa grupa Liego.
Postać macierzy zgadza się z otrzymaną wcześniej w
Przykładzie~\ref{prz:bloch}; tu wykorzystujemy ją jako
globalny opis rozmaitości i działań grupowych.
\end{przyklad}

\section{Od wektora przy jedności do algebry Liego}
\label{sec:algebra-liego}

\begin{definicja}[Pole lewostronnie niezmiennicze i algebra Liego]
\label{def:algebra-liego}
\index{pole lewostronnie niezmiennicze}
\index{algebra Liego!definicja}
Pole $V$ na $G$ jest \emph{lewostronnie niezmiennicze}, gdy
$(L_g)_*V=V$ dla każdego $g\in G$. Dla $X\in T_eG$ oznaczamy
\[
 X^L_g:=(\dd L_g)_eX.
\]
\emph{Algebrą Liego} grupy $G$ nazywamy przestrzeń
$\mathfrak g=T_eG$ z nawiasem
\[
 [X,Y]_{\mathfrak g}:=[X^L,Y^L]_e.
\]
Wektor przy jedności opisuje ruch nieskończenie mały; translacje
przenoszą ten sam ruch do wszystkich punktów. Przykładami są
$\mathfrak{gl}(n,\R)=M_n(\R)$ i, jak pokażemy niżej,
$\mathfrak t^k\cong\R^k$ z zerowym nawiasem.
\end{definicja}

\begin{twierdzenie}[Niezmienniczość i zamknięcie na nawias]
\label{thm:lie-algebra-fields}
Odwzorowanie $X\mapsto X^L$ jest izomorfizmem liniowym
$T_eG$ na przestrzeń pól lewostronnie niezmienniczych.
Nawias dwóch takich pól jest znowu lewostronnie niezmienniczy.
W szczególności $\mathfrak g$ jest algebrą Liego: jej nawias jest
dwuliniowy, antysymetryczny i spełnia tożsamość Jacobiego.
\end{twierdzenie}
\begin{proof}
Gładkość $X^L$ wynika z gładkości różniczki mnożenia po
ograniczeniu do $G\times\{e\}$. Dla $g,h\in G$ mamy
$L_g\circ L_h=L_{gh}$, więc reguła łańcucha daje
\[
 (\dd L_g)_h X^L_h
 = \bigl((\dd L_g)_h\circ(\dd L_h)_e\bigr)X
 = (\dd L_{gh})_eX=X^L_{gh}.
\]
Jest to niezmienniczość. Jeśli $V$ jest dowolnym polem
lewostronnie niezmienniczym, to po podstawieniu $h=e$
otrzymujemy $V_g=(\dd L_g)_eV_e$; zatem $V=(V_e)^L$.
Wartość przy $e$ jest więc odwrotnością odwzorowania
$X\mapsto X^L$, a oba odwzorowania są liniowe.

Dla dyfeomorfizmu $F$ zachodzi
$F_*[V,W]=[F_*V,F_*W]$. Można to sprawdzić bezpośrednio
na funkcji $f$: definicja pola przeniesionego daje
$(F_*V)(f)=(V(f\circ F))\circ F^{-1}$, a po dwukrotnym
zastosowaniu tego wzoru różnica złożeń staje się
$([V,W](f\circ F))\circ F^{-1}$.
Dokładniej, w drugim zastosowaniu skraca się
$F^{-1}\circ F$, toteż
$(F_*V)((F_*W)f)=(V(W(f\circ F)))\circ F^{-1}$.
Po zamianie $V,W$ i odjęciu otrzymujemy wskazany nawias.
Stosując tę równość do $F=L_g$ i $V=X^L$, $W=Y^L$,
dostajemy
\[
 (L_g)_*[X^L,Y^L]
 =[(L_g)_*X^L,(L_g)_*Y^L]=[X^L,Y^L].
\]
Własności nawiasu pól zostały dowiedzione w
Fakcie~\ref{fakt:nawias-wlasnosci}. Przez izomorfizm
$X\mapsto X^L$ przechodzą na nawias w $T_eG$.
\end{proof}

\begin{figure}[!htbp]
\centering
\begin{tikzpicture}[>=Stealth,line cap=round,line join=round]
 \fill[manifoldblue!7] (0,0) circle(1.6);
 \draw[manifoldblue!75!black,thick] (0,0) circle(1.6);
 \draw[manifoldblue!30,densely dashed] (0,0)--(1.6,0) (0,0)--(55:1.6);
 \draw[->,manifoldgold!85!black,thick]
  (6:1.6) arc[start angle=6,end angle=48,radius=1.6];
 \fill[manifoldred] (1.6,0) circle(2.4pt) node[below right=2pt] {$e=1$};
 \fill[manifoldred] (55:1.6) circle(2.4pt) node[above right=2pt] {$g=e^{i\theta}$};
 \draw[->,very thick,manifoldred] (1.6,0)--(1.6,.85)
  node[right,font=\small] {$X=i$};
 \draw[->,very thick,manifoldred] (55:1.6)--++(145:.85)
  node[above left,font=\small] {$X^L_g=gi$};
 \node[font=\small,align=left,anchor=west,text width=6.3cm]
  at (3.1,.6)
  {$L_g$ obraca okrąg o kąt $\theta$.
   Jego różniczka obraca również wektor styczny:
   $(\dd L_g)_e(i)=gi$.\\[4pt]
   Z jednego wektora przy $e$ otrzymujemy całe pole.
   Wektor $gi$ jest styczny do okręgu w $g$.};
\end{tikzpicture}
\caption{Pole lewostronnie niezmiennicze na $S^1$.
Niebieski dysk jest tylko tłem rysunku; grupą jest jego brzeg.
Pole powstaje przez różniczkowanie translacji, bez wyboru metryki.}
\label{fig:lie-pole-okrag}
\end{figure}

\begin{przyklad}[Komutator macierzy]
\label{prz:lie-bracket-matrix}
\index{komutator macierzy}
Na otwartym $\mathrm{GL}(n,\R)$ pole odpowiadające macierzy
$X$ ma wartość $X^L_A=AX$. Dla $Y^L_A=AY$ różniczka
odwzorowania $A\mapsto AY$ w kierunku $H$ jest równa $HY$.
Wzór współrzędny na nawias z~\eqref{eq:nawias-wspolrzedne}
daje więc
\[
 [X^L,Y^L]_A
 =\D(Y^L)_A[AX]-\D(X^L)_A[AY]
 =AXY-AYX=A(XY-YX).
\]
Stąd $[X,Y]_{\mathfrak{gl}}=XY-YX$. To samo obliczenie,
po ograniczeniu pól stycznych do podgrupy macierzy, obowiązuje
dla $\mathrm O(n)$, $\mathrm{SO}(n)$ i $\mathrm{SU}(2)$.
W szczególności
$\mathfrak{so}(n)=\{X:X^T=-X\}$ jest zamknięta na komutator,
bo $(XY-YX)^T=Y^TX^T-X^TY^T=YX-XY=-(XY-YX)$.

Dla $\mathrm{SU}(2)$ różniczkowanie $U^*U=I$ i $\det U=1$
w punkcie $I$ daje odpowiednio $X^*+X=0$ i
$\operatorname{tr}X=0$; różniczka wyznacznika w kierunku
$X$ wynosi $\operatorname{tr}X$, co widać z rozwinięcia
$\det(I+tX)=1+t\operatorname{tr}X+t^2\det X$.
Ponieważ obie strony mają wymiar
rzeczywisty trzy, otrzymujemy
\[
 \mathfrak{su}(2)
 =\left\{\begin{pmatrix}i\alpha&z\\-\bar z&-i\alpha\end{pmatrix}
 :\alpha\in\R,\ z\in\C\right\}.
\]
Przestrzeń po prawej ma trzy rzeczywiste współrzędne:
$\alpha,\Re z,\Im z$. Inkluzja wynika z różniczkowania warunków,
a równość wymiarów z $\mathrm{SU}(2)\cong S^3$,
więc nie zakładamy bez dowodu, że każde rozwiązanie
zlinearyzowanych równań jest styczne.
Na torusie translacje w mapach ilorazowych mają stałą
różniczkę równą identyczności. Odpowiadające im pola mają
stałe współczynniki, więc ich nawias znika:
$\mathfrak t^k\cong(\R^k,0)$.
\end{przyklad}

\begin{przyklad}[Grupa Heisenberga i geometria Nil]
\label{prz:lie-heisenberg}
\index{grupa Heisenberga}
Współrzędnym $(x,y,z)\in\R^3$ przypiszmy macierz
$H(x,y,z)=\begin{psmallmatrix}1&x&z\\0&1&y\\0&0&1\end{psmallmatrix}$.
W iloczynie $H(x,y,z)H(x',y',z')$ współczynnik $(1,3)$
to $1\cdot z'+x\cdot y'+z\cdot1$, a współczynniki $(1,2)$
i $(2,3)$ są sumami. Stąd prawo grupowe ma postać
\[
 (x,y,z)(x',y',z')=(x+x',y+y',z+z'+xy').
\]
Element neutralny to $(0,0,0)$, a rozwiązanie równania
$(x,y,z)(x',y',z')=(0,0,0)$ daje
$(x,y,z)^{-1}=(-x,-y,-z+xy)$.
Po podstawieniu w odwrotnej kolejności także dostajemy jedność.
Łączność dziedziczymy z macierzy. Oba działania są wielomianowe
we współrzędnych, więc jest to grupa Liego na $\R^3$.

Dla ustalonego $g=(x,y,z)$ różniczkujemy
$L_g(x',y',z')=(x+x',y+y',z+z'+xy')$
w $(x',y',z')=0$. Macierz różniczki to
$\begin{psmallmatrix}1&0&0\\0&1&0\\0&x&1\end{psmallmatrix}$.
Jej kolumny dają pola odpowiadające bazie w jedności:
$E_1=\partial_x$, $E_2=\partial_y+x\partial_z$,
$E_3=\partial_z$.
We wzorze~\eqref{eq:nawias-wspolrzedne} jedyną niezerową
pochodną współczynnika jest $\partial_x(x)=1$, zatem
$[E_1,E_2]=E_3$, $[E_1,E_3]=[E_2,E_3]=0$.

Równania przepływów możemy rozwiązać od razu:
$\Phi_a^{E_1}(x,y,z)=(x+a,y,z)$,
$\Phi_b^{E_2}(x,y,z)=(x,y+b,z+xb)$,
$\Phi_c^{E_3}(x,y,z)=(x,y,z+c)$.
W drugim wzorze $x$ jest stałe w czasie, więc równania
$y'=1$, $z'=x$ dają przyrosty $b$ i $xb$.
Zaczynając w jedności i wykonując kolejno przepływy
$E_1,E_2,-E_1,-E_2$ przez czasy $a,b,a,b$, otrzymujemy
\[
 (0,0,0)\longmapsto(a,0,0)\longmapsto(a,b,ab)
 \longmapsto(0,b,ab)\longmapsto(0,0,ab).
\]
Rzut na płaszczyznę $(x,y)$ jest zamkniętym prostokątem,
ale w grupie pozostaje przyrost $ab$ w kierunku $E_3$.
Tak nawias $[E_1,E_2]=E_3$ mierzy nieprzemienność ruchów.

Ramkę dualną w sensie Definicji~\ref{def:ramka-lokalna}
tworzą formy $\dd x,\dd y,\dd z-x\dd y$.
Metryka $g_N=(\dd x)^2+(\dd y)^2+(\dd z-x\dd y)^2$
z Przykładu~\ref{prz:thur-dop-metryki-nil-sol}
jest lewostronnie niezmiennicza. Aby to sprawdzić, dla
$g_0=(x_0,y_0,z_0)$ oznaczmy współrzędne $L_{g_0}$
przez $(\widetilde x,\widetilde y,\widetilde z)$.
Mamy $\dd\widetilde x=\dd x$, $\dd\widetilde y=\dd y$ oraz
$\dd\widetilde z=\dd z+x_0\dd y$, więc
$\dd\widetilde z-\widetilde x\dd\widetilde y
=\dd z+x_0\dd y-(x_0+x)\dd y=\dd z-x\dd y$.
Cofnięcie zachowuje każdy z trzech składników metryki.
\end{przyklad}

\begin{figure}[!htbp]
\centering
\begin{tikzpicture}[>=Stealth,line cap=round,line join=round]
 \coordinate (O) at (0,0);
 \coordinate (P) at (2,0);
 \coordinate (Q) at (2.8,3.4);
 \coordinate (R) at (.8,3.4);
 \coordinate (S) at (0,2.8);
 \fill[manifoldblue!7] (0,0)--(2,0)--(2.8,.6)--(.8,.6)--cycle;
 \draw[manifoldblue!35,densely dashed] (P)--(2.8,.6)--(.8,.6)--(O);
 \draw[manifoldblue!30,densely dashed] (2.8,.6)--(Q) (.8,.6)--(R);
 \draw[->,very thick,manifoldred] (O)--(P) node[midway,below=3pt] {$E_1$};
 \draw[->,very thick,manifoldgreen!70!black] (P)--(Q)
  node[midway,right=3pt] {$E_2$};
 \draw[->,very thick,manifoldred] (Q)--(R) node[midway,above=3pt] {$-E_1$};
 \draw[->,very thick,manifoldgreen!70!black] (R)--(S);
 \node[manifoldgreen!70!black] at (.1,3.68) {$-E_2$};
 \draw[<->,manifoldgold!85!black,thick] (-.3,0)--(-.3,2.8)
  node[midway,left=3pt] {$ab$};
 \foreach \p in {O,P,Q,R,S}{\fill[manifoldblue!75!black] (\p) circle(2pt);}
 \node[below left=3pt] at (O) {$e$};
 \node[left=12pt] at (S) {$(0,0,ab)$};
 \node[font=\small,align=left,anchor=west,text width=5.7cm] at (4.2,1.8)
  {Cztery ruchy zamykają prostokąt w rzucie $(x,y)$.
   Przyrost wysokości pozostaje równy $ab$.\\[4pt]
   Odcinki mają czasy $a,b,a,b$, gdzie $a,b>0$.};
\end{tikzpicture}
\caption{Nieprzemienność w grupie Heisenberga.
Rysunek jest rzutem ukośnym współrzędnych $(x,y,z)$.
Przerywany prostokąt leży na $z=0$; wzdłuż boku $E_2$
na poziomie $x=a$ rośnie także współrzędna $z$.}
\label{fig:lie-heisenberg-komutator}
\end{figure}

\section{Podgrupy jednoparametrowe i eksponenta}
\label{sec:lie-exponential}

\begin{definicja}[Podgrupa jednoparametrowa]
\label{def:one-parameter}
\index{podgrupa jednoparametrowa}
\emph{Podgrupą jednoparametrową} nazywamy gładki
homomorfizm $\gamma:(\R,+)\to G$.
Warunek homomorfizmu ma postać
$\gamma(s+t)=\gamma(s)\gamma(t)$ i $\gamma(0)=e$.
Na $S^1$ przykładem jest $t\mapsto e^{it}$; ruch wraca
do jedności po czasie $2\pi$, więc homomorfizm nie musi być
różnowartościowy.
\end{definicja}

\begin{twierdzenie}[Wektor styczny wyznacza jedną podgrupę]
\label{thm:one-param-lie}
Dla każdego $X\in\mathfrak g$ istnieje dokładnie jedna
podgrupa jednoparametrowa $\gamma_X$ taka, że
$\gamma_X'(0)=X$. Jest ona krzywą całkową pola $X^L$
przez $e$ i zachodzi
\[
 \Phi_t^{X^L}(g)=g\gamma_X(t)
 \qquad (g\in G,\ t\in\R).
\]
\end{twierdzenie}
\begin{proof}
Twierdzenie~\ref{thm:przeplyw} daje lokalną krzywą
całkową $\gamma$ pola $X^L$ z $\gamma(0)=e$.
Jeżeli $\eta(t)$ jest krzywą całkową przez $h$, to
$(L_{h^{-1}}\circ\eta)(t)$ również jest krzywą całkową $X^L$:
niezmienniczość pola i reguła łańcucha przenoszą równanie
$\eta'=X^L_\eta$ na tę krzywą. Dokładniej,
\[
 (L_{h^{-1}}\circ\eta)'(t)
 =(\dd L_{h^{-1}})_{\eta(t)}X^L_{\eta(t)}
 =X^L_{h^{-1}\eta(t)}.
\]
Jej wartością w zerze jest $e$. Jednoznaczność rozwiązań
z Rozdziału~\ref{ch:pola} daje zatem, dla małych $t$,
$\Phi_t^{X^L}(h)=h\gamma(t)$.

Wykażmy, że pole jest zupełne. Wybierzmy $\varepsilon>0$,
tak aby $\gamma(t)$ istniała dla $|t|<\varepsilon$.
Z niezmienniczości $X^L$ krzywa
$s\mapsto h\gamma(s)$, $|s|<\varepsilon$, jest
rozwiązaniem przez \emph{dowolny} $h\in G$ przez taki sam
przedział czasu. Gdyby maksymalna krzywa przez $e$ miała
skończony prawy koniec $b$, wybierając $t_0<b$ z
$b-t_0<\varepsilon/2$, przedłużylibyśmy ją wzorem
$\gamma(t_0+s)=\gamma(t_0)\gamma(s)$ do
$t_0+s>b$, wbrew maksymalności.
Na wspólnym przedziale stara i nowa krzywa zgadzają się:
po zmianie czasu na $s=0$ rozwiązują to samo równanie
z tą samą wartością $\gamma(t_0)$. Jednoznaczność
zapewnia więc także gładkie sklejenie. Lewy koniec traktujemy
tak samo. Stąd $\gamma$ istnieje dla wszystkich $t$.

Dla ustalonego $s$ krzywa $t\mapsto\gamma(s)\gamma(t)$
jest całkowa i w chwili $t=0$ przechodzi przez
$\gamma(s)$. Krzywa $t\mapsto\gamma(s+t)$ ma te same
własności. Jednoznaczność daje
$\gamma(s+t)=\gamma(s)\gamma(t)$, zatem $\gamma$ jest
podgrupą jednoparametrową. Ponownie przez translację
jej przepływ przez $g$ ma podaną postać.

Jeśli $\beta:\R\to G$ jest inną gładką podgrupą
jednoparametrową z $\beta'(0)=X$, to
\[
 \beta'(t)=\left.\frac{\dd}{\dd s}\right|_{s=0}
 \beta(t)\beta(s)=(\dd L_{\beta(t)})_eX=X^L_{\beta(t)}.
\]
Jednoznaczność rozwiązania z warunkiem początkowym w $e$
daje $\beta=\gamma$.
\end{proof}

Lewostronna niezmienniczość pola i prawa translacja w jego
przepływie są zgodne. Dla ustalonego $g$ różniczkujemy
krzywą $g\gamma_X(t)$ po czasie, a nie po punkcie $g$.
Niezmienniczość przenosi równanie z $\gamma_X(t)$ do
$g\gamma_X(t)$ przez $L_g$. Natomiast dla ustalonego czasu
$t$ całe odwzorowanie $g\mapsto g\gamma_X(t)$ jest
$R_{\gamma_X(t)}$.

\begin{definicja}[Odwzorowanie wykładnicze]
\label{def:lie-exp}
\index{odwzorowanie wykładnicze!grupy Liego}
\emph{Odwzorowaniem wykładniczym} grupy Liego nazywamy
\[
 \exp_G:\mathfrak g\longrightarrow G,\qquad
 \exp_G(X)=\gamma_X(1).
\]
Wektor $X$ określa prędkość przy jedności, a $\exp_G(X)$
jest punktem osiągniętym po czasie jeden przy tej stałej,
przenoszonej translacjami prędkości. Dla torusa z
Przykładu~\ref{prz:torus-liego}, przy identyfikacji
$T_eT^k\cong\R^k$ z mapy ilorazowej, mamy
$\exp_{T^k}(v)=[v]$.
\end{definicja}

Jest to eksponenta \emph{grupowa}. Jej konstrukcja używa
mnożenia i pola $X^L$, bez metryki ani koneksji.
Odwzorowanie $\exp_p$ z Definicji~\ref{def:exp-geodezyjna}
powstawało natomiast z równania geodezyjnych.
Nie utożsamiamy tych konstrukcji bez dodatkowego uzasadnienia.

\begin{stwierdzenie}[Własności eksponenty]
\label{prop:lie-exp-properties}
Odwzorowanie $\exp_G$ jest gładkie,
$(\dd\exp_G)_0=\id_{\mathfrak g}$ oraz
\[
 \exp_G((s+t)X)=\exp_G(sX)\exp_G(tX),\qquad
 \exp_G(-X)=\exp_G(X)^{-1}.
\]
W szczególności $\exp_G$ jest dyfeomorfizmem pewnego
otoczenia zera na otoczenie $e$.
\end{stwierdzenie}
\begin{proof}
Liniowa zależność $X\mapsto X^L$ i gładkość różniczki
mnożenia sprawiają, że $(X,g)\mapsto X^L_g$ jest
gładkie. Na rozmaitości $\mathfrak g\times G$ rozważmy
pole $(X,g)\mapsto(0,X^L_g)$. Jego przepływ ma postać
$(X,g)\mapsto(X,\Phi_t^{X^L}(g))$ i jest określony
dla wszystkich czasów na mocy
Twierdzenia~\ref{thm:one-param-lie}. Gładkość przepływu
z Twierdzenia~\ref{thm:przeplyw} daje więc gładkość
$(X,t)\mapsto\gamma_X(t)$, a w szczególności $\exp_G$.
Jednoznaczność podgrupy o zadanej pochodnej
daje $\gamma_X(t)=\exp_G(tX)$: dla każdej liczby $a$
krzywa $s\mapsto\gamma_X(as)$ ma pochodną $aX$ w zerze,
więc $\gamma_{aX}(1)=\gamma_X(a)$.
Prawo grupowe z poprzedniego dowodu daje pierwszą równość,
a podstawienie $s=-t$ daje drugą.

Dla $X\in\mathfrak g$ mamy
\[
 (\dd\exp_G)_0(X)
 =\left.\frac{\dd}{\dd t}\right|_{t=0}\exp_G(tX)
 =\gamma_X'(0)=X.
\]
Różniczka jest zatem izomorfizmem. Twierdzenie o funkcji
odwrotnej z Rozdziału~\ref{ch:rozmaitosci} daje ostatni wniosek.
\end{proof}

\begin{przyklad}[Eksponenta macierzowa]
\label{prz:matrix-exp-lie}
\index{eksponenta macierzowa}
Dla $X\in M_n(\R)$ lub $M_n(\C)$ szereg
\[
 E_X(t)=\sum_{j=0}^{\infty}\frac{t^jX^j}{j!}
\]
jest zbieżny bezwzględnie i jednostajnie z pochodnymi
na każdym zwartym przedziale $t$: wybieramy normę
operatorową macierzy. Dla $|t|\leq T$ norma składnika nie
przekracza $T^j\|X\|^j/j!$, a norma jego $m$-tej
pochodnej dla $j\geq m$ nie przekracza
$\|X\|^m(T\|X\|)^{j-m}/(j-m)!$. Szeregi tych oszacowań
są zbieżne. Różniczkowanie wyraz po wyrazie i zmiana
indeksu $\ell=j-1$ dają
\[
 E_X'(t)=\sum_{j=1}^\infty\frac{t^{j-1}X^j}{(j-1)!}
 =\left(\sum_{\ell=0}^\infty\frac{t^\ell X^\ell}{\ell!}\right)X
 =E_X(t)X,\qquad E_X(0)=I.
\]
Macierz $X$ komutuje z każdym $E_X(t)$, więc
reguła iloczynu i reguła łańcucha dają
\[
 \frac{\dd}{\dd t}\bigl(E_X(t)E_X(-t)\bigr)
 =E_X(t)X E_X(-t)-E_X(t)E_X(-t)X=0.
\]
Ostatnia równość używa komutowania $X$ z $E_X(-t)$.
Wartość iloczynu
w zerze to $I$. Zatem $E_X(t)^{-1}=E_X(-t)$
i cała krzywa leży w $\mathrm{GL}(n,\mathbb K)$.
Spełnia równanie pola $X^L_A=AX$ na
$\mathrm{GL}(n,\mathbb K)$; z jednoznaczności wynika
\[
 \exp_{\mathrm{GL}(n,\mathbb K)}(X)
 =E_X(1)=\sum_{j=0}^{\infty}\frac{X^j}{j!}.
\]
Jeżeli $X^T=-X$, to
$E_X(t)^T=E_{-X}(t)=E_X(t)^{-1}$;
transpozycja przechodzi pod znak szeregu,
a $(X^j)^T=(X^T)^j=(-X)^j$.
Ciągłość wyznacznika od $t=0$ pokazuje
$E_X(t)\in\mathrm{SO}(n)$.
Jeśli
$X=\begin{psmallmatrix}i\alpha&z\\-\bar z&-i\alpha\end{psmallmatrix}
\in\mathfrak{su}(2)$, połóżmy $r=\sqrt{\alpha^2+|z|^2}$.
Mnożenie macierzy daje
\[
 X^2=\begin{pmatrix}
 -\alpha^2-z\bar z&i\alpha z-i\alpha z\\
 -i\alpha\bar z+i\alpha\bar z&-\bar z z-\alpha^2
 \end{pmatrix}=-r^2I.
\]
Wobec tego $X^{2j}=(-r^2)^jI$ i
$X^{2j+1}=(-r^2)^jX$.
Parzyste wyrazy szeregu sumują się do $\cos(rt)I$,
a nieparzyste do $\sin(rt)X/r$, czyli
$E_X(t)=\cos(rt)I+\sin(rt)X/r$ dla $r>0$.
Przy $r=0$ mamy $X=0$, więc $E_X(t)=I$.
Jeśli oznaczymy $c=\cos(rt)$, $s=\sin(rt)$, to
\[
 E_X(t)=
 \begin{pmatrix}c+i\alpha s/r&zs/r\\
 -\bar z s/r&c-i\alpha s/r\end{pmatrix},
 \qquad
 \det E_X(t)=c^2+\frac{\alpha^2+|z|^2}{r^2}s^2=c^2+s^2.
\]
Ponadto $X^*=-X$, zatem
$E_X(t)^*E_X(t)=(cI-sX/r)(cI+sX/r)
=c^2I-s^2X^2/r^2=(c^2+s^2)I=I$.
Otrzymaliśmy unitarność i wyznacznik $1$, czyli
$E_X(t)\in\mathrm{SU}(2)$, bez korzystania z ogólnego
wzoru na wyznacznik eksponenty.

Dla $J=\begin{psmallmatrix}0&-1\\1&0\end{psmallmatrix}$
zachodzi $J^2=-I$, więc rozdzielenie parzystych i
nieparzystych potęg daje
\[
 \exp(\theta J)=
 \begin{pmatrix}\cos\theta&-\sin\theta\\
 \sin\theta&\cos\theta\end{pmatrix}.
\]
Torus w modelu $(S^1)^k$ ma z kolei eksponentę
$(\theta_1,\ldots,\theta_k)\mapsto
(e^{i\theta_1},\ldots,e^{i\theta_k})$ i jądro
$(2\pi\Z)^k$; współrzędne $\theta$ są tutaj
$2\pi$ razy współrzędnymi $v$ z definicji ilorazowej.
Lokalna odwracalność eksponenty nie oznacza więc
globalnej różnowartościowości. W $\mathrm{GL}(n,\R)$
eksponenta nie pokrywa nawet całej grupy: każdy
$\exp X$ leży na drodze $t\mapsto\exp(tX)$ od $I$,
więc ma dodatni wyznacznik.

Przeszkoda występuje także pośród macierzy połączonych drogą
z $I$. Dla $A=\operatorname{diag}(-1,-2)$ mamy $\det A=2>0$.
Gdyby $A=e^X$ dla rzeczywistej macierzy $X$, to $XA=AX$,
bo $X$ komutuje ze swoim szeregiem wykładniczym.
Pisząc $X=\begin{psmallmatrix}p&q\\r&s\end{psmallmatrix}$,
porównujemy współczynniki $(1,2)$ i $(2,1)$:
$-2q=-q$ i $-r=-2r$, czyli $q=r=0$.
Wówczas $e^X=\operatorname{diag}(e^p,e^s)$ ma dodatnie
wyrazy diagonalne, sprzeczność.
Macierz $A$ należy jednak do składowej jedności:
droga $t\mapsto\exp(\pi tJ)\operatorname{diag}(1,2^t)$,
$0\leq t\leq1$, zaczyna się w $I$, kończy w $A$
i ma przez cały czas wyznacznik $2^t>0$.

Dla grupy Heisenberga z Przykładu~\ref{prz:lie-heisenberg}
wektor $(a,b,c)$ przy jedności reprezentuje macierz
$N=\begin{psmallmatrix}0&a&c\\0&0&b\\0&0&0\end{psmallmatrix}$.
Jedynym niezerowym współczynnikiem $N^2$ może być
$(N^2)_{13}=ab$, a $N^3=0$. Szereg kończy się więc po trzech wyrazach:
\[
 \exp(tN)=I+tN+\tfrac12t^2N^2
 =H\bigl(ta,tb,tc+\tfrac12t^2ab\bigr).
\]
W szczególności $\exp(a,b,c)=(a,b,c+ab/2)$ w tych
współrzędnych grupowych. Przy $X=(1,0,0)$, $Y=(0,1,0)$
mamy $\exp(X+Y)=(1,1,1/2)$, natomiast prawo mnożenia daje
$\exp(X)\exp(Y)=(1,1,1)$.
Własność z~\ref{prop:lie-exp-properties} dotyczy jednej
prostej $\R X$; nie uprawnia do takiej równości dla dowolnych $X,Y$.
\end{przyklad}

\begin{przyklad}[Zamknięta i gęsta podgrupa na torusie]
\label{prz:lie-torus-podgrupy}
Na $T^2=\R^2/\Z^2$ pole odpowiadające $v\in\R^2$ ma
stałe współczynniki, więc jego krzywa przez $0$ to
$\gamma_v(t)=[tv]$.
Powrót do jedności zachodzi dokładnie wtedy, gdy
$tv\in\Z^2$.
Dla $v=(1,1/2)$ oznacza to $t\in2\Z$.
Krzywa schodzi więc do injektywnego odwzorowania
$\R/2\Z\to T^2$ z niezerową różniczką.
Jest zanurzeniem okręgu: ciągła bijekcja na obraz
ze zwartej dziedziny do przestrzeni Hausdorffa jest
homeomorfizmem, a niezerowa różniczka daje immersję.
Jej obraz jest zwarty, zatem domknięty.

Dla $v=(1,\alpha)$, gdzie $\alpha\notin\Q$,
warunki $t\in\Z$, $\alpha t\in\Z$ wymuszają $t=0$.
Zatem $\gamma_v$ jest injektywna i ma wszędzie
niezerową różniczkę, ale jej obraz jest gęsty w torusie.
Uzasadnijmy gęstość bez dodatkowego twierdzenia.
Oznaczmy przez $\{r\}=r-\lfloor r\rfloor$ część ułamkową.
Wśród $N+1$ różnych liczb
$\{0\alpha\},\{1\alpha\},\ldots,\{N\alpha\}$
dwie należą do tego samego z $N$ przedziałów długości $1/N$.
Ich dodatnia różnica ma postać
$\delta=m\alpha-k$, gdzie $m,k\in\Z$ i $0<\delta<1/N$.
Dla dowolnego $r\in[0,1)$ liczba $j=\lfloor r/\delta\rfloor$
spełnia $0\leq r-j\delta<\delta$ i $j\delta<1$.
Ponieważ $jm\alpha=jk+j\delta$, klasa $[jm\alpha]$
przybliża $[r]$ z błędem mniejszym niż $1/N$.
Klasy $[n\alpha]$, $n\in\Z$, są więc gęste w $\R/\Z$.
Dla ustalonego $s$ punkty
$\gamma_v(s+n)=([s],[\alpha s+n\alpha])$
wypełniają gęsto cały okrąg o pierwszej współrzędnej $[s]$.
To dowodzi gęstości obrazu w $T^2$.

Ta injektywna immersja $\R\to T^2$ nie jest zanurzeniem.
Powyższa konstrukcja daje $m_N\in\Z$ z
$0<\delta_N=m_N\alpha-k_N<1/N$.
Musi zachodzić $|m_N|\to\infty$: dla ograniczonego
zbioru niezerowych $m_N$ odległości $m_N\alpha$ od $\Z$
miałyby dodatnie minimum.
Tymczasem $\gamma_v(m_N)=([0],[\delta_N])\to([0],[0])$,
a $m_N$ nie dąży do $0$. Odwrotność na obrazie
nie jest więc ciągła w jedności.
Podgrupa jednoparametrowa jest z definicji homomorfizmem;
jej obraz nie musi być zanurzoną podrozmaitością.
\end{przyklad}

\begin{figure}[!htbp]
\centering
\begin{tikzpicture}[>=Stealth,line cap=round,line join=round,
 flow/.style={manifoldred,thick,postaction={decorate},
 decoration={markings,mark=at position .55 with {\arrow{Stealth}}}}]
 \foreach \shift in {0,5.3}{
  \begin{scope}[xshift=\shift cm]
   \fill[manifoldblue!7] (0,0) rectangle (3.4,3.4);
   \draw[manifoldblue!75!black,thick] (0,0) rectangle (3.4,3.4);
   \draw[->,manifoldblue!60] (.45,-.12)--(.95,-.12);
   \draw[->,manifoldblue!60] (.45,3.52)--(.95,3.52);
   \draw[->,manifoldblue!60] (-.12,.45)--(-.12,.95);
   \draw[->,manifoldblue!60] (3.52,.45)--(3.52,.95);
   \fill[manifoldred] (0,0) circle(2.5pt);
   \node[below left=4pt,font=\small] at (0,0) {$e$};
  \end{scope}
 }
 \draw[flow] (0,0)--(3.4,1.7);
 \draw[flow] (0,1.7)--(3.4,3.4);
 \node at (1.7,3.95) {$v=(1,1/2)$};
 \node[font=\small] at (1.7,-.65) {$0\leq t\leq2$: jeden okres};
 \begin{scope}[xshift=5.3cm]
  \draw[flow] (0.0000,0.0000)--(2.4042,3.4000);
  \draw[manifoldred,thick] (2.4042,0.0000)--(3.4000,1.4083);
  \draw[manifoldred,thick] (0.0000,1.4083)--(1.4083,3.4000);
  \draw[manifoldred,thick] (1.4083,0.0000)--(3.4000,2.8167);
  \draw[flow] (0.0000,2.8167)--(0.4125,3.4000);
  \draw[manifoldred,thick] (0.4125,0.0000)--(2.8167,3.4000);
  \draw[manifoldred,thick] (2.8167,0.0000)--(3.4000,0.8250);
  \draw[manifoldred,thick] (0.0000,0.8250)--(1.8208,3.4000);
  \draw[flow] (1.8208,0.0000)--(3.4000,2.2333);
  \draw[manifoldred,thick] (0.0000,2.2333)--(0.8250,3.4000);
  \draw[manifoldred,thick] (0.8250,0.0000)--(3.2291,3.4000);
  \draw[manifoldred,thick] (3.2291,0.0000)--(3.4000,0.2416);
  \draw[flow] (0.0000,0.2416)--(2.2333,3.4000);
  \draw[manifoldred,thick] (2.2333,0.0000)--(3.4000,1.6500);
 \end{scope}
 \node at (7,3.95) {$v=(1,\sqrt2)$};
 \node[font=\small] at (7,-.65) {$0\leq t\leq6$: fragment krzywej};
\end{tikzpicture}
\caption{Podgrupy jednoparametrowe $t\mapsto[tv]$
w kwadratach fundamentalnych torusa. Przeciwne boki są
utożsamione zgodnie z niebieskimi strzałkami.
Przy nachyleniu $1/2$ dwa odcinki zamykają krzywą;
przy $\sqrt2$ kolejne przejścia dają gęstą orbitę,
co uzasadniono w Przykładzie~\ref{prz:lie-torus-podgrupy}.}
\label{fig:lie-torus-podgrupy}
\end{figure}

\section{Sprzężenie i działanie adjungowane}
\label{sec:lie-adjoint}

\begin{definicja}[Działanie adjungowane]
\label{def:lie-ad}
\index{działanie adjungowane}
Dla $g\in G$ \emph{sprzężenie} przez $g$ to automorfizm
$c_g:G\to G$, $c_g(h)=ghg^{-1}$. Jego różniczkę przy $e$
oznaczamy
\[
 \operatorname{Ad}_g:=(\dd c_g)_e:
 \mathfrak g\longrightarrow\mathfrak g.
\]
Ponieważ $c_g\circ c_h=c_{gh}$, otrzymujemy
$\operatorname{Ad}_{gh}=\operatorname{Ad}_g\circ\operatorname{Ad}_h$.
Jest to liniowy sposób zapisania, jak zmiana punktu odniesienia
przez sprzężenie zmienia prędkości przy jedności.
W grupie macierzy $\operatorname{Ad}_gX=gXg^{-1}$,
co wynika przez zróżniczkowanie
$g\exp_G(tX)g^{-1}$ w zerze.
\end{definicja}

\begin{twierdzenie}[Różniczka działania adjungowanego]
\label{thm:lie-ad-bracket}
Odwzorowanie $g\mapsto\operatorname{Ad}_g$ jest gładkim
homomorfizmem $G\to\mathrm{GL}(\mathfrak g)$.
Dla $X,Y\in\mathfrak g$ zachodzi
\begin{equation}\label{eq:lie-ad-bracket}
 \left.\frac{\dd}{\dd t}\right|_{t=0}
 \operatorname{Ad}_{\exp_G(tX)}Y=[X,Y]_{\mathfrak g}.
\end{equation}
Ponadto $\operatorname{Ad}_g[X,Y]_{\mathfrak g}
=[\operatorname{Ad}_gX,\operatorname{Ad}_gY]_{\mathfrak g}$.
\end{twierdzenie}
\begin{proof}
Gładkość sprzężenia jako odwzorowania $(g,h)\mapsto ghg^{-1}$
oraz gładkość jego różniczki dają gładkość $g\mapsto
\operatorname{Ad}_g$. Reguła łańcucha i
$c_g\circ c_h=c_{gh}$ dają własność homomorfizmu
i odwrotność $\operatorname{Ad}_{g^{-1}}$.
Ponieważ $c_g\circ L_h=L_{c_g(h)}\circ c_g$,
różniczkowanie w $e$ daje
\[
 \begin{aligned}
 ((c_g)_*Y^L)_{c_g(h)}
 &=(\dd c_g)_h\bigl((\dd L_h)_eY\bigr)\\
 &=(\dd L_{c_g(h)})_e\bigl((\dd c_g)_eY\bigr)
   =(\operatorname{Ad}_gY)^L_{c_g(h)}.
 \end{aligned}
\]
Punkt $c_g(h)$ przebiega całą grupę, toteż otrzymujemy równość pól:
\begin{equation}\label{eq:ad-pole}
 (c_g)_*Y^L=(\operatorname{Ad}_gY)^L.
\end{equation}

Ustalmy $X$ i połóżmy $a_t=\exp_G(tX)$.
Rodzina $c_{a_t}$ jest przepływem pola
$V=X^R-X^L$, gdzie $X^R_h=(\dd R_h)_eX$.
Istotnie, różniczkowanie $a_t h a_{-t}$ przy $t=0$
daje $(\dd R_h)_eX-(\dd L_h)_eX$.
Możemy rozdzielić te dwa wkłady przez odwzorowanie
$F(s,t)=a_s h a_{-t}$: pochodna $F(t,t)$ w zerze
jest sumą $\partial_sF(0,0)+\partial_tF(0,0)$.
Pierwsza pochodna pochodzi z krzywej $a_s h$,
druga z $h a_{-t}$ i wnosi minus z reguły łańcucha.
Ponadto $c_{a_s}\circ c_{a_t}=c_{a_{s+t}}$.
Różniczkując tę tożsamość względem $s$ w zerze,
otrzymujemy równanie przepływu $V$ w każdym punkcie
krzywej $t\mapsto c_{a_t}(h)$.
Przepływem $X^R$ jest $h\mapsto a_t h$, a przepływem
$Y^L$ jest $h\mapsto h\exp_G(tY)$. Translacje lewe
i prawe komutują, bo $a(hb)=(ah)b$, zatem z
Stwierdzenia~\ref{prop:inwarianty} wynika
$[X^R,Y^L]=0$. Korzystając z konwencji pochodnej
Liego z Twierdzenia~\ref{thm:lie-nawias}, mamy
$(c_{a_t})_*Y^L=(c_{a_{-t}})^*Y^L$.
Pochodna cofnięcia przez przepływ o czasie $\tau$ to
$\Lie_VY^L=[V,Y^L]$; tutaj $\tau=-t$, więc
reguła łańcucha wprowadza znak minus. Dlatego
\[
 \left.\frac{\dd}{\dd t}\right|_{t=0}
 (c_{a_t})_*Y^L
 =-[X^R-X^L,Y^L]=[X^L,Y^L].
\]
Ocena przy $e$ i~\eqref{eq:ad-pole} dają~\eqref{eq:lie-ad-bracket}.
Na koniec używamy naturalności nawiasu względem $c_g$
z dowodu Twierdzenia~\ref{thm:lie-algebra-fields}.
Ponieważ $c_g(e)=e$, ocena przy jedności daje
\[
 \begin{aligned}
 \operatorname{Ad}_g[X,Y]_{\mathfrak g}
 &=((c_g)_*[X^L,Y^L])_e\\
 &=[(c_g)_*X^L,(c_g)_*Y^L]_e
 =[\operatorname{Ad}_gX,\operatorname{Ad}_gY]_{\mathfrak g},
 \end{aligned}
\]
gdzie w ostatnim kroku podstawiliśmy~\eqref{eq:ad-pole}.
\end{proof}

\begin{przyklad}[Macierze i torus]
\label{prz:ad-matrix-torus}
Dla grupy macierzy możemy niezależnie sprawdzić znak
w~\eqref{eq:lie-ad-bracket}:
\[
 \left.\frac{\dd}{\dd t}\right|_0
 e^{tX}Ye^{-tX}=XY-YX.
\]
Pierwszy czynnik ma w zerze pochodną $X$,
ostatni $-X$, a wartości obu wynoszą $I$.
Reguła iloczynu daje więc $XYI+IY(-X)$.
Dla konkretnego przykładu weźmy
$X=\begin{psmallmatrix}1&0\\0&0\end{psmallmatrix}$,
$Y=\begin{psmallmatrix}0&1\\0&0\end{psmallmatrix}$.
Ponieważ $e^{tX}=\operatorname{diag}(e^t,1)$,
bezpośrednio mnożymy
\[
 \operatorname{Ad}_{e^{tX}}Y
 =\begin{pmatrix}e^t&0\\0&1\end{pmatrix}
  \begin{pmatrix}0&1\\0&0\end{pmatrix}
  \begin{pmatrix}e^{-t}&0\\0&1\end{pmatrix}
 =e^tY.
\]
Pochodna w zerze to $Y$, zgodnie z
$XY=Y$, $YX=0$ i $[X,Y]=Y$.
W grupie przemiennej każde sprzężenie jest identycznością.
Stąd $\operatorname{Ad}$ torusa jest trywialne, a jej
różniczka i nawias algebry Liego znikają.
W teorii Cherna--Weila z Rozdziału~\ref{ch:chern-weil}
warunek niezmienniczości wielomianu oznacza właśnie
niezmienniczość względem takich przekształceń
$\operatorname{Ad}_g$.
\end{przyklad}

\begin{stwierdzenie}[Różniczka homomorfizmu i naturalność eksponenty]
\label{prop:lie-homomorfizm}
Niech $F:G\to H$ będzie gładkim homomorfizmem grup Liego.
Wówczas $A=(\dd F)_e:\mathfrak g\to\mathfrak h$
jest homomorfizmem algebr Liego, czyli
$A[X,Y]_{\mathfrak g}=[AX,AY]_{\mathfrak h}$.
Ponadto $F(\exp_G(tX))=\exp_H(tAX)$ dla każdego $t\in\R$.
\end{stwierdzenie}
\begin{proof}
Homomorfizm spełnia $F\circ L_g=L_{F(g)}\circ F$.
Różniczkując w $e$, otrzymujemy
$(\dd F)_gX^L_g=(AX)^L_{F(g)}$.
Nie zakładamy, że $F$ jest dyfeomorfizmem:
jest to równość wektorów po zastosowaniu różniczki.
Dla każdej funkcji gładkiej $f$ na $H$ znaczy ona
$X^L(f\circ F)=((AX)^Lf)\circ F$.
Stosujemy ją dwukrotnie, najpierw do $Y$, potem do $X$,
i odejmujemy wzór z zamienionymi $X,Y$:
\[
 [X^L,Y^L](f\circ F)
 =\bigl([(AX)^L,(AY)^L]f\bigr)\circ F.
\]
Ocena w $e$ i definicja nawiasu algebry Liego dają
$A[X,Y]=[AX,AY]$, ponieważ równość zachodzi na każdej $f$.

Złożenie $F\circ\gamma_X$ jest gładkim homomorfizmem
$\R\to H$, a jego pochodna w zerze to $AX$.
Jednoznaczność z Twierdzenia~\ref{thm:one-param-lie}
daje $F\circ\gamma_X=\gamma_{AX}$, czyli drugą tezę.
\end{proof}

Oznaczmy przez $\operatorname{ad}_X$ endomorfizm
$Y\mapsto[X,Y]$ przestrzeni $\mathfrak g$.
Wzór~\eqref{eq:lie-ad-bracket} mówi, że
$(\dd\operatorname{Ad})_e(X)=\operatorname{ad}_X$.
Stosując Stwierdzenie~\ref{prop:lie-homomorfizm}
do $F=\operatorname{Ad}:G\to\mathrm{GL}(\mathfrak g)$, otrzymujemy
\[
 \operatorname{Ad}_{\exp_G(tX)}
 =\exp_{\mathrm{GL}(\mathfrak g)}(t\operatorname{ad}_X).
\]
Po wyborze bazy prawa strona jest szeregiem macierzowym
z Przykładu~\ref{prz:matrix-exp-lie}. Na wektorze $Y$
ma wartość $Y+t[X,Y]+\tfrac12t^2[X,[X,Y]]+\cdots$:
każda kolejna potęga $\operatorname{ad}_X$ dodaje jeden nawias z $X$.

\section{Od SU(2) do obrotów przestrzeni}
\label{sec:su2-so3-lie}

W Przykładzie~\ref{prz:spin3} grupa
$\operatorname{Spin}(3)$ została utożsamiona z jednostkowymi
kwaternionami. Tutaj połączymy ten opis z macierzami
$\mathrm{SU}(2)$ i obliczymy samo nakrycie.

Zapiszmy kwaternion jako $q=a+b\mathbf j$, gdzie $a,b\in\C$,
$\mathbf j^2=-1$ i $\mathbf jz=\bar z\mathbf j$.
Iloczyn ma wtedy postać
\[
 (a+b\mathbf j)(c+d\mathbf j)
 =(ac-b\bar d)+(ad+b\bar c)\mathbf j.
\]
Istotnie, po rozdzieleniu nawiasów dostajemy $ac+ad\mathbf j
+b\mathbf jc+b\mathbf jd\mathbf j$.
W trzecim wyrazie używamy $\mathbf jc=\bar c\mathbf j$,
a w czwartym $\mathbf jd\mathbf j=\bar d\mathbf j^2=-\bar d$.
Odwzorowanie $q\mapsto U(a,b)$ z~\eqref{eq:su2-macierz}
zachowuje ten iloczyn, co widać po pomnożeniu macierzy:
jej pierwszy wiersz ma współczynniki $ac-b\bar d$
i $ad+b\bar c$. Ponadto
$|q|^2=|a|^2+|b|^2$. Jest to więc gładki izomorfizm
grup jednostkowych kwaternionów z $\mathrm{SU}(2)$.
Odwrotność odczytuje pierwszy wiersz macierzy, więc także
jest gładka. Ten sam wzór $U(a,b)$ zachowuje mnożenie
również bez warunku jednostkowej normy.

\begin{twierdzenie}[Jawne dwukrotne nakrycie]
\label{thm:su2-so3-cover-lie}
\index{nakrycie!$\mathrm{SU}(2)$ nad $\mathrm{SO}(3)$}
Niech $\Im\mathbb H$ będzie przestrzenią
kwaternionów o zerowej części rzeczywistej, utożsamioną
z $\R^3$ przez bazę $(\mathbf i,\mathbf j,\mathbf k)$.
Wtedy
\[
 \rho:\mathrm{SU}(2)\longrightarrow\mathrm{SO}(3),
 \qquad \rho(q)(v)=qvq^{-1},\quad
 v\in\Im\mathbb H,
\]
jest gładkim surjektywnym homomorfizmem o jądrze
$\{1,-1\}$. Każdy punkt $\mathrm{SO}(3)$ ma otoczenie
równomiernie pokryte przez dwa otwarte zbiory w
$\mathrm{SU}(2)$; zatem $\rho$ jest dwukrotnym nakryciem.
\end{twierdzenie}
\begin{proof}
\emph{Krok 1: dlaczego sprzężenie jest obrotem?}
Najpierw uzasadnijmy własności iloczynu.
Ze wzoru $\det U(q)=|q|^2$ i mnożenia macierzy wynika
$|qr|^2=\det(U(q)U(r))=|q|^2|r|^2$,
więc $|qr|=|q||r|$.
Sprzężenie kwaternionowe ma postać
$\overline{a+b\mathbf j}=\bar a-b\mathbf j$.
Podstawienie do wzoru na iloczyn daje
$q\bar q=\bar q q=|q|^2$, zatem $q^{-1}=\bar q$ dla $|q|=1$.
Jeśli $p=p_0+\sum_i p_i\mathbf e_i$,
$r=r_0+\sum_i r_i\mathbf e_i$ w bazie
$\mathbf e_1=\mathbf i$, $\mathbf e_2=\mathbf j$,
$\mathbf e_3=\mathbf k$, to
$\Re(pr)=p_0r_0-\sum_i p_ir_i=\Re(rp)$:
kwadraty generatorów to $-1$, a mieszane iloczyny są urojone.
Stąd
$\Re(qvq^{-1})=\Re(q^{-1}(qv))=\Re v$
oraz $|qvq^{-1}|=|v|$.
Odwzorowanie $v\mapsto qvq^{-1}$ jest rzeczywiście liniowe.
Zachowuje więc także iloczyn skalarny, otrzymany z normy przez
$2\langle v,w\rangle=|v+w|^2-|v|^2-|w|^2$.
Ogranicza się zatem do izometrii $\Im\mathbb H$.
Zależność od $q$ jest
gładka, bo $q^{-1}=\bar q$ na sferze jednostkowej.
Grupa $S^3$ jest spójna drogowo: każdy jej element
można zapisać jako $\cos t+u\sin t$ dla jednostkowego
urojonego $u$ (przy $q=\pm1$ wybór $u$ jest dowolny).
Dla $q=a+w$ z $w\ne0$ wybieramy $t=\arccos a\in(0,\pi)$
i $u=w/|w|$, gdyż $|w|=\sqrt{1-a^2}=\sin t$.
Droga $\tau\mapsto\cos(\tau t)+u\sin(\tau t)$,
$0\leq\tau\leq1$, łączy $1$ z $q$ w sferze jednostkowej.
Wyznacznik izometrii $\rho(q)$ jest więc stale równy
$+1$, jak dla $q=1$. Łączność mnożenia daje
$\rho(qr)=\rho(q)\circ\rho(r)$:
po podstawieniu dowolnego $v$ obie strony dają
$q(rv r^{-1})q^{-1}$.

\emph{Krok 2: wyprowadzenie kąta obrotu.}
Dla urojonych kwaternionów $u,v$ mamy
$uv=-\langle u,v\rangle+u\times v$.
W części rzeczywistej pozostaje $-\sum_i u_iv_i$.
Współczynnik przy $\mathbf i$ to $u_2v_3-u_3v_2$,
bo $\mathbf j\mathbf k=\mathbf i$ i
$\mathbf k\mathbf j=-\mathbf i$; dwie pozostałe
współrzędne powstają przez cykliczną zamianę indeksów.
To dokładnie współrzędne iloczynu wektorowego.
Niech $|u|=1$ i
$q=\cos(\theta/2)+u\sin(\theta/2)$.
Ponieważ $uv-vu=2u\times v$ i
$uvu=v-2\langle u,v\rangle u$, rozwinięcie
$(c+su)v(c-su)$ dla
$c=\cos(\theta/2)$, $s=\sin(\theta/2)$ daje
\begin{align*}
 \rho(q)v
 &=c^2v+cs(uv-vu)-s^2uvu\\
 &=\cos\theta\,v+
   \sin\theta\,(u\times v)+
   (1-\cos\theta)\langle u,v\rangle u.
\end{align*}
Druga użyta tożsamość wynika z
$uv+vu=-2\langle u,v\rangle$: mnożymy ją z prawej strony
przez $u$ i podstawiamy $u^2=-1$.
Przejście między dwoma wierszami rachunku wykorzystuje
$c^2-s^2=\cos\theta$, $2cs=\sin\theta$ i
$2s^2=1-\cos\theta$.
Dla $v=\lambda u+w$, $w\perp u$, wzór upraszcza się do
$\rho(q)v=\lambda u+\cos\theta\,w+\sin\theta\,(u\times w)$.
Oś $u$ pozostaje nieruchoma.
Jeśli $w$ jest jednostkowy, to $(w,u\times w)$ jest
dodatnią bazą ortonormalną płaszczyzny $u^\perp$.
Podstawienie kolejno $w$ i $u\times w$, z użyciem
$u\times(u\times w)=-w$, daje macierz
$\begin{psmallmatrix}\cos\theta&-\sin\theta\\
\sin\theta&\cos\theta\end{psmallmatrix}$ na tej płaszczyźnie.
Jest to zatem obrót o kąt $\theta$ wokół osi $u$.

\emph{Krok 3: surjektywność.}
Każda macierz $A\in\mathrm{SO}(3)$ ma oś stałą:
\[
 \det(I-A)=\det(I-A^T)
 =\det(A^T)\det(A-I)
 =\det(A-I)=-\det(I-A),
\]
więc $\det(I-A)=0$. Wybierzmy jednostkowy
$u\in\ker(I-A)$.
W rachunku użyliśmy $I-A^T=A^T(A-I)$ oraz
$\det(A-I)=(-1)^3\det(I-A)$.
Ponieważ $A^Tu=u$, dla $w\perp u$ mamy
$\langle Aw,u\rangle=\langle w,A^Tu\rangle=0$.
Płaszczyzna $u^\perp$ jest więc niezmiennicza.
W rozkładzie $\R^3=\R u\oplus u^\perp$ macierz ma bloki
$1$ i $B$, zatem $\det B=\det A=1$.
Na $u^\perp$ macierz $A$ jest izometrią o wyznaczniku $+1$.
Orientujemy tę płaszczyznę
tak, by $(u,\text{dodatnia baza }u^\perp)$ była dodatnia
w $\R^3$; wtedy $A|_{u^\perp}$ jest obrotem pewnego
kąta $\theta$. Aby to zobaczyć, zapiszmy pierwszą kolumnę $B$
jako $(a,b)^T$ z $a^2+b^2=1$.
Jedyna prostopadła kolumna jednostkowa dająca dodatni wyznacznik
to $(-b,a)^T$. Wybór $\cos\theta=a$, $\sin\theta=b$
daje właśnie macierz obrotu.
Wzór z Kroku 2 daje kwaternion,
który realizuje dokładnie $A$; $\rho$ jest surjektywna.

\emph{Krok 4: jądro i włókna.}
Jeśli $\rho(q)$ jest identycznością, to $q$ komutuje
ze wszystkimi urojonymi $v$. Po zapisaniu $q=a+u$
z $a\in\R$, $u\in\Im\mathbb H$,
otrzymujemy $0=qv-vq=2u\times v$ dla każdego $v$.
Stąd $u=0$ i, przez $|q|=1$, $q=\pm1$.
Rzeczywiście, dla $u\ne0$ wybór niezerowego $v\perp u$
dałby $|u\times v|=|u||v|>0$, sprzeczność.
Oba kwaterniony $\pm1$ istotnie działają identycznie.
Jeśli $\rho(r)=\rho(q)$, to $\rho(q^{-1}r)=I$,
czyli $r=\pm q$. Każde włókno ma zatem dokładnie dwa punkty.

\emph{Krok 5: lokalne arkusze nakrycia.}
Dla $u\in\Im\mathbb H$
krzywa $q(t)=\cos(t|u|)+(u/|u|)\sin(t|u|)$
(dla $u=0$ stała) ma pochodną $u$ w zerze.
Różniczkując sprzężenie, dostajemy
\[
 (\dd\rho)_1(u)(v)
 =\left.\frac{\dd}{\dd t}\right|_{t=0}q(t)vq(t)^{-1}
 =uv-vu=2u\times v.
\]
Pochodna $q(t)^{-1}$ w zerze jest równa $-u$:
różniczkowanie $q(t)q(t)^{-1}=1$ daje $u+(q^{-1})'(0)=0$.
W bazie $(\mathbf i,\mathbf j,\mathbf k)$, dla
$u=(u_1,u_2,u_3)$, otrzymujemy macierz
\[
 (\dd\rho)_1(u)=
 2\begin{pmatrix}0&-u_3&u_2\\u_3&0&-u_1\\-u_2&u_1&0\end{pmatrix}.
\]
Jest to izomorfizm $\Im\mathbb H\to
\mathfrak{so}(3)$: jest różnowartościowy, bo
$u\times v=0$ dla wszystkich $v$ wymusza $u=0$,
a obie przestrzenie mają wymiar trzy.
Ponieważ $\rho\circ L_q=L_{\rho(q)}\circ\rho$,
różniczkowanie i odwrócenie różniczki translacji dają
\[
 (\dd\rho)_q=(\dd L_{\rho(q)})_I\circ
             (\dd\rho)_1\circ(\dd L_{q^{-1}})_q.
\]
Wszystkie trzy czynniki są izomorfizmami.
Twierdzenie
o funkcji odwrotnej daje więc otoczenie $W$ jedności,
na którym $\rho$ jest dyfeomorfizmem. Zmniejszmy $W$
tak, aby $\Re q>0$ na $W$; wtedy $W\cap(-W)=\varnothing$. Połóżmy
$V=\rho(W)$. Każdy punkt $V$ ma dokładnie
dwa przeciwobrazy, zatem
$\rho^{-1}(V)=W\sqcup(-W)$ i oba ograniczenia są
dyfeomorfizmami na $V$: na $-W$ odwrotność otrzymujemy
z odwrotności na $W$ przez pomnożenie przez $-1$.
Dla dowolnego obrotu $A_0=\rho(q_0)$ otoczenie $A_0V$
ma arkusze $q_0W$ i $q_0(-W)$.
Równość $\rho(q_0w)=A_0\rho(w)$ przenosi na nie
oba lokalne dyfeomorfizmy. To jest
definicja dwukrotnego nakrycia z~\ref{def:nakrycie}.
\end{proof}

\begin{figure}[!htbp]
\centering
\begin{tikzpicture}[>=Stealth,line cap=round,line join=round]
 \fill[manifoldblue!7] (0,0) circle(1.55);
 \draw[manifoldblue!75!black,thick] (0,0) circle(1.55);
 \draw[manifoldblue!30,densely dashed] (215:1.55)--(35:1.55);
 \draw[manifoldblue!30,densely dashed] (0,0)--(1.55,0);
 \draw[->,manifoldgold!85!black,thick]
  (.65,0) arc[start angle=0,end angle=35,radius=.65];
 \node[font=\small] at (18:.94) {$\varphi$};
 \fill[manifoldred] (35:1.55) circle(2.7pt) node[above right=2pt] {$q$};
 \fill[manifoldred] (215:1.55) circle(2.7pt) node[below left=2pt] {$-q$};
 \fill[manifoldblue!75!black] (1.55,0) circle(2pt) node[below right=2pt] {$1$};
 \fill[manifoldblue!75!black] (-1.55,0) circle(2pt) node[above left=2pt] {$-1$};
 \node[font=\small] at (0,2.25) {$q=\cos\varphi+\mathbf k\sin\varphi\in S^3$};
 \draw[->,very thick,manifoldblue!75!black] (2.15,0)--(4.25,0)
  node[midway,above=4pt] {$\rho$}
  node[midway,below=5pt,font=\small] {$q,\,-q$};
 \begin{scope}[xshift=6.4cm]
  \fill[manifoldblue!7] (0,0) circle(1.55);
  \draw[manifoldblue!75!black,thick] (0,0) circle(1.55);
  \draw[manifoldblue!30,densely dashed] (0,0)--(1.55,0) (0,0)--(70:1.55);
  \draw[->,manifoldgold!85!black,thick]
   (.65,0) arc[start angle=0,end angle=70,radius=.65];
  \node[font=\small] at (35:1.03) {$2\varphi$};
  \fill[manifoldred] (70:1.55) circle(2.7pt) node[above right=2pt] {$R_{2\varphi}$};
  \fill[manifoldblue!75!black] (1.55,0) circle(2pt) node[below right=2pt] {$I$};
  \node[font=\small] at (0,2.25) {obroty wokół osi $\mathbf k$};
 \end{scope}
 \node[font=\small,align=center] at (3.2,-2.25)
  {$\varphi:0\longrightarrow\pi$ daje $q:1\longrightarrow-1$,\\
   a kąt obrotu zmienia się od $0$ do $2\pi$.};
\end{tikzpicture}
\caption{Podwojenie kąta w nakryciu $\mathrm{SU}(2)\to\mathrm{SO}(3)$.
Po lewej pokazano jeden okrąg wielki w $S^3$, po prawej okrąg
podgrupy obrotów wokół ustalonej osi.
Punkty antypodalne $q$ i $-q$ dają ten sam obrót:
zmiana $\varphi$ o $\pi$ zmienia kąt $2\varphi$ o $2\pi$.}
\label{fig:lie-su2-podwojenie-kata}
\end{figure}

\begin{przyklad}[Obrót o kąt prosty]
\label{prz:lie-obrot-kwaternion}
Weźmy oś $u=\mathbf k$ i kąt $\theta=\pi/2$.
Połowa kąta to $\pi/4$, zatem wzór z dowodu
Twierdzenia~\ref{thm:su2-so3-cover-lie} daje
$q=(1+\mathbf k)/\sqrt2$.
Ponieważ $\mathbf k=\mathbf i\mathbf j$, w zapisie $a+b\mathbf j$
mamy $a=1/\sqrt2$, $b=i/\sqrt2$.
Odpowiadająca macierz i obrót to
\[
 U(q)=\frac1{\sqrt2}\begin{pmatrix}1&i\\i&1\end{pmatrix},
 \qquad
 \rho(q)=\begin{pmatrix}0&-1&0\\1&0&0\\0&0&1\end{pmatrix}.
\]
Drugą macierz otrzymujemy, podstawiając do wzoru obrotu
kolejno $\mathbf i,\mathbf j,\mathbf k$:
$\mathbf k\times\mathbf i=\mathbf j$,
$\mathbf k\times\mathbf j=-\mathbf i$,
a oś $\mathbf k$ pozostaje nieruchoma.
Są to kolejne kolumny macierzy.
Kwaternion $-q$ daje te same kolumny, bo w
$(-q)v(-q)^{-1}$ oba minusy się znoszą.
Natomiast $q^2=(1+2\mathbf k+\mathbf k^2)/2=\mathbf k$,
$q^4=\mathbf k^2=-1$ i $q^8=1$: cztery kolejne obroty o $\pi/2$
wracają do $I$ w $\mathrm{SO}(3)$, a ich podniesienie
kończy się w $-1$. Dopiero osiem takich kroków wraca do $1$.
\end{przyklad}

W szczególności $\mathrm{SU}(2)\cong\operatorname{Spin}(3)$
jako grupy Liego, a powyższe $\rho$ jest tym samym
nakryciem co w Twierdzeniu~\ref{thm:spin-nakrycie}.
Rzeczywiście w algebrze Clifforda z
Przykładu~\ref{prz:spin3} element
$\omega=e_1e_2e_3$ jest centralny, a
$\mathcal J(v)=-v\omega$ wysyła bazę $(e_1,e_2,e_3)$ na
$(I,J,K)$ przestrzeni urojonych kwaternionów.
Przesunięcie $e_i$ przez pozostałe dwa generatory wnosi
dwa minusy, więc $e_i\omega=\omega e_i$.
Na przykład $-e_1\omega=-e_1^2e_2e_3=e_2e_3=I$;
cykliczna zamiana indeksów daje $J,K$.
Ponieważ $\omega$ komutuje z $s\in\operatorname{Spin}(3)$,
otrzymujemy $\mathcal J(svs^{-1})=s\mathcal J(v)s^{-1}$, czyli obie
konstrukcje nakrycia są zgodne pod tym utożsamieniem.
Jego różniczka zachowuje nawias na mocy
Stwierdzenia~\ref{prop:lie-homomorfizm} i jest bijekcją,
jak obliczyliśmy w Kroku 5. Jest więc izomorfizmem algebr Liego
$\mathfrak{su}(2)\cong\mathfrak{so}(3)$, lecz same grupy
Liego nie są izomorficzne. Istotnie,
$\mathrm{SU}(2)\cong S^3$ jest jednospójna na mocy
Przykładu~\ref{prz:pi1-sfera}. Droga
$q(t)=\cos(\pi t)+\mathbf i\sin(\pi t)$, $0\leq t\leq1$,
łączy $1$ z $-1$, a jej obraz przez $\rho$ jest pętlą
w $\mathrm{SO}(3)$. Gdyby ta pętla była ściągalna,
Twierdzenie~\ref{thm:podnoszenie-nakrycie} pozwoliłoby
podnieść jej homotopię do homotopii dróg o stałych
końcach. Początek podniesienia jest stale równy $1$.
Jego koniec zależy ciągle od parametru homotopii i należy
do dyskretnego zbioru $\{1,-1\}$, więc pozostaje równy $-1$.
Tymczasem podniesienie końcowej pętli stałej, zaczynające się
w $1$, jest stale równe $1$. To sprzeczność.
Zatem $\pi_1(\mathrm{SO}(3))$ nie jest trywialna.
Możemy obliczyć ją dokładnie przez
Twierdzenie~\ref{thm:deck-normalizator}.
Nakrycie ma jednospójną dziedzinę, więc podgrupa $H$
w tym twierdzeniu jest trywialna, a grupa transformacji
nakrycia jest izomorficzna z $\pi_1(\mathrm{SO}(3))$.
Każda transformacja wysyła $q$ na $q$ albo $-q$.
Wybór znaku jest lokalnie stały, zatem stały na spójnej $S^3$.
Obie możliwości są transformacjami, więc otrzymujemy
$\pi_1(\mathrm{SO}(3))\cong\Z/2\Z$.

\par\smallskip
{\footnotesize\noindent\emph{Dalsza lektura:}
\href{https://archive.org/download/theotransformation01liesrich/theotransformation01liesrich.pdf}
{S.\ Lie i F.\ Engel, \emph{Theorie der Transformationsgruppen},
t.~I (1888)}, oryginalne ujęcie grup przekształceń i pól
infinitesymalnych;
\href{https://books.google.com/books/about/Theory_of_Lie_Groups_I.html?id=bm4GAQAAIAAJ}
{C.\ Chevalley, \emph{Theory of Lie Groups}, t.~I (1946)},
systematyczne ujęcie grup gładkich i ich algebr.
Współczesny wykład z dowodami:
\href{https://www.staff.science.uu.nl/~ban00101/lecnotes/lie2010.pdf}
{E.\ P.\ van den Ban, \emph{Lie Groups} (2010)},
§§2--4 (pola, eksponenta i $\operatorname{Ad}$), §7 (podgrupy)
oraz §10 ($\mathrm{SU}(2)$ i $\mathrm{SO}(3)$).\par}
